% =====================================================================
%  preamble.tex  --  SHARED PREAMBLE FOR CS-4031 COMPILER CONSTRUCTION
%  Fall 2026, FAST-NU Lahore
% =====================================================================

% ---------- Page geometry ----------
\usepackage[a4paper, margin=2.2cm, headsep=14pt, footskip=28pt]{geometry}

% ---------- Fonts / encoding ----------
\usepackage[utf8]{inputenc}
\usepackage[T1]{fontenc}

% ---------- Math ----------
\usepackage{amsmath, amssymb, amsthm}

% ---------- Colors ----------
\usepackage[table]{xcolor}
\definecolor{nuNavy}{HTML}{102A43}      % primary heading color
\definecolor{nuBlue}{HTML}{1B5FA8}      % accent / links
\definecolor{nuLight}{HTML}{EAF2FB}     % light panel background
\definecolor{dbGreen}{HTML}{1F7A4D}     % "Dragon Book / examinable" boxes
\definecolor{dbGreenBg}{HTML}{EAF7EF}
\definecolor{enrichGold}{HTML}{9A6A00}  % "enrichment / beyond DB" boxes
\definecolor{enrichGoldBg}{HTML}{FFF6E0}
\definecolor{examRed}{HTML}{B3261E}     % "exam focus" boxes
\definecolor{examRedBg}{HTML}{FCEBEA}
\definecolor{histPurple}{HTML}{5B3A8E}
\definecolor{histPurpleBg}{HTML}{F2EBFB}
\definecolor{takeawayTeal}{HTML}{0F6F6A}
\definecolor{takeawayTealBg}{HTML}{E7F6F5}
\definecolor{lightgray}{HTML}{F5F5F5}

% ---------- Headers / footers ----------
\usepackage{fancyhdr}
\pagestyle{fancy}
\fancyhf{}
\renewcommand{\headrulewidth}{0.6pt}
\renewcommand{\footrulewidth}{0.3pt}
\fancyhead[L]{\small\sffamily\color{nuNavy} CS-4031 Compiler Construction}
\fancyhead[R]{\small\sffamily\color{nuNavy} \rightmark}
\fancyfoot[L]{\small\sffamily\color{gray} FAST-NU Lahore, Fall 2026}
\fancyfoot[C]{\small\sffamily\color{gray} \thepage}
\fancyfoot[R]{\small\sffamily\color{gray} Notes by course staff}

% ---------- Section styling ----------
\usepackage{titlesec}
\titleformat{\section}
  {\normalfont\Large\bfseries\color{nuNavy}}
  {\thesection}{0.6em}{}
\titleformat{\subsection}
  {\normalfont\large\bfseries\color{nuBlue}}
  {\thesubsection}{0.6em}{}
\titlespacing*{\section}{0pt}{1.4em}{0.8em}
\titlespacing*{\subsection}{0pt}{1.1em}{0.5em}

% ---------- Lists, tables, misc ----------
\usepackage{enumitem}
\setlist{topsep=4pt, itemsep=2pt, parsep=0pt}
\usepackage{array, booktabs, tabularx, longtable, multirow}
\usepackage{graphicx}
\usepackage{tikz}
\usetikzlibrary{arrows.meta, positioning, shapes.geometric, fit, backgrounds, calc}
\usepackage{float}
\usepackage[colorlinks=true, linkcolor=nuBlue, citecolor=nuBlue, urlcolor=nuBlue]{hyperref}

% ---------- Table-of-contents number width ----------
% Default article.cls reserves just enough horizontal space in the TOC for a
% single-digit section number. Once any lecture grows past 9 top-level
% sections (as a merged/longer lecture can), two-digit numbers like "3.10"
% collide visually with the title text right after them ("3.10Extensions").
% tocloft lets us widen that reserved space once, here, so every lecture
% is protected regardless of how many sections it ends up with.
\usepackage{tocloft}
\setlength{\cftsecnumwidth}{2.8em}
\setlength{\cftsubsecnumwidth}{3.4em}
\renewcommand{\theHsection}{L\thelecnum.\arabic{section}}
\renewcommand{\theHsubsection}{L\thelecnum.\arabic{section}.\arabic{subsection}}
\renewcommand{\theHsubsubsection}{L\thelecnum.\arabic{section}.\arabic{subsection}.\arabic{subsubsection}}

% ---------- Cross-reference machinery (no hard-coded section/lecture numbers) ----------
% Every lecture resets \setcounter{section}{0} at its own top (see \lecturemeta / the
% \setcounter{section}{0} call in each lecture file), so within a single lecture, sections
% would normally print as a bare "1", "2", "3", ... -- indistinguishable at a glance from the
% *next* lecture's own "1", "2", "3", ... once everything is stitched into one combined PDF.
% Prefixing the *displayed* section number with the lecture number fixes this without touching
% any individual lecture file: "3.2" unambiguously means "Lecture 3, Section 2" everywhere,
% including in the combined table of contents.
\renewcommand{\thesection}{\thelecnum.\arabic{section}}

% \Sref{<label>} is the standard way to refer to a section from prose, in place of typing a
% literal number. Pair every \section/\subsection that is ever referred to elsewhere with a
% \label{sec:...} right after it, then write \Sref{sec:...} wherever the old text would have
% hard-coded "\S2.4" or similar. If the numbering ever shifts (a new section inserted, content
% reordered), every \Sref pointing at that label updates automatically on the next compile.
\newcommand{\Sref}[1]{\S\ref{#1}}


% ---------- Box framework (tcolorbox) ----------
\usepackage[most]{tcolorbox}
\tcbuselibrary{skins, breakable}

\newtcolorbox{lecbox}[3]{
  breakable, enhanced, sharp corners=downhill,
  colback=#2, colframe=#1,
  boxrule=0.6pt, left=6pt, right=6pt, top=4pt, bottom=4pt,
  fonttitle=\bfseries\sffamily, coltitle=white,
  colbacktitle=#1,
  title={#3},
  attach boxed title to top left={yshift=-0.3mm, xshift=2mm},
  boxed title style={sharp corners, colframe=#1, boxrule=0pt},
  top=8mm
}

\newenvironment{coreidea}[1][Core Idea (Dragon Book)]
  {\begin{lecbox}{dbGreen}{dbGreenBg}{#1}}
  {\end{lecbox}}

\newenvironment{examfocus}[1][Exam Focus]
  {\begin{lecbox}{examRed}{examRedBg}{#1}}
  {\end{lecbox}}

\newenvironment{enrichment}[1][Beyond the Dragon Book -- Enrichment]
  {\begin{lecbox}{enrichGold}{enrichGoldBg}{#1}}
  {\end{lecbox}}

\newenvironment{historynote}[1][A Bit of History]
  {\begin{lecbox}{histPurple}{histPurpleBg}{#1}}
  {\end{lecbox}}

\newenvironment{takeaways}[1][Key Takeaways]
  {\begin{lecbox}{takeawayTeal}{takeawayTealBg}{#1}}
  {\end{lecbox}}

\newtcolorbox{defbox}[1]{
  breakable, enhanced, colback=nuLight, colframe=nuBlue,
  boxrule=0.5pt, left=6pt, right=6pt, top=3pt, bottom=3pt,
  fonttitle=\bfseries, coltitle=nuNavy,
  title={Definition: #1}
}
\newtcolorbox{examplebox}[1]{
  breakable, enhanced, colback=white, colframe=gray!60,
  boxrule=0.5pt, left=6pt, right=6pt, top=3pt, bottom=3pt,
  fonttitle=\bfseries, coltitle=black,
  title={Example: #1}
}

% ---------- Lecture metadata machinery ----------
\newcommand{\thelecnum}{}
\newcommand{\lecshorttitle}{}
\newcommand{\lecclo}{}
\newcommand{\lecdbrefs}{}
\newcommand{\lecotherrefs}{}
\newcommand{\lecfulltitle}{}

\newcommand{\lecturemeta}[5]{%
  \global\renewcommand{\thelecnum}{#1}%
  \global\renewcommand{\lecfulltitle}{#2}%
  \global\renewcommand{\lecclo}{#3}%
  \global\renewcommand{\lecdbrefs}{#4}%
  \global\renewcommand{\lecotherrefs}{#5}%
}
\newcommand{\lecshort}[1]{%
  \global\renewcommand{\lecshorttitle}{#1}%
  \markright{Lecture \thelecnum\ --- #1}%
}

% ---------- Per-lecture \label and table-of-contents entry ----------
% \lecturelabel defines \label{lec:<N>} so other lectures can write "Lecture~\ref{lec:5}"
% instead of a hard-coded "Lecture 5" that silently goes stale if lectures are ever
% renumbered or reordered. \label normally captures whatever \thesection/\thechapter/etc.
% last set as \@currentlabel; since a bare lecture title isn't a sectioning command, we set
% \@currentlabel to \thelecnum ourselves first, so \ref{lec:<N>} prints just the number.
% \makeatletter/\makeatother are required here: \@currentlabel contains "@", which is not a
% letter by default in a plain \input file like this one (only inside .cls/.sty loading does
% LaTeX make that automatic) -- without this, "\@currentlabel" tokenizes as the unrelated
% control-symbol "\@" followed by the literal word "currentlabel", silently missing the
% kernel's real \@currentlabel entirely and leaving every \ref{lec:...} blank.
\makeatletter
\newcommand{\lecturelabel}{%
  \global\edef\@currentlabel{\thelecnum}%
  \label{lec:\thelecnum}%
}
\makeatother

% \lecturetocentry inserts one bold, page-level line into the *combined* table of contents
% for this lecture, at the same outline depth LaTeX normally reserves for \part -- one level
% above \section. This is what turns the combined TOC from a flat, ambiguous run of
% "1, 2, 3, 1, 2, 3, 4, ..." into a clearly hierarchical "Lecture 1: <title> / 1.1 ... / 1.2
% ... // Lecture 2: <title> / 2.1 ... / ...". It only affects the TOC listing -- it does not
% print a page break or heading in the body, since \makelecturetitle already provides that.
\newcommand{\lecturetocentry}{%
  \phantomsection
  \addcontentsline{toc}{part}{Lecture \thelecnum: \lecshorttitle}%
}

\newcommand{\makelecturetitle}{%
  \lecturelabel
  \lecturetocentry
  \begin{center}
    {\small\sffamily\color{gray} THE NATIONAL UNIVERSITY OF COMPUTER AND EMERGING SCIENCES, LAHORE}\\[2pt]
    {\small\sffamily\color{gray} CS-4031 Compiler Construction --- Fall 2026}\\[10pt]
    {\Huge\bfseries\color{nuNavy} Lecture \thelecnum}\\[4pt]
    {\LARGE\bfseries\color{nuBlue} \lecfulltitle}
  \end{center}
  \vspace{6pt}
  \noindent
  \begin{tcolorbox}[breakable, enhanced, colback=nuLight, colframe=nuNavy,
      boxrule=0.6pt, left=8pt, right=8pt, top=4pt, bottom=4pt]
    \small
    \begin{tabularx}{\textwidth}{@{}l X@{}}
      \textbf{CLO(s) covered:} & \lecclo \\
      \textbf{Dragon Book [DB]:} & \lecdbrefs \\
      \textbf{Other references:} & \lecotherrefs \\
    \end{tabularx}
  \end{tcolorbox}
  \vspace{4pt}
}

\newcommand{\beyonddb}{\textcolor{enrichGold}{\small\textbf{[beyond DB]}}\ }
\newcommand{\dbtag}{\textcolor{dbGreen}{\small\textbf{[DB]}}\ }

% ---------- Bibliography convention ----------
% Each lecture file ends its own "References for This Lecture" section with a
% \begin{thebibliography}{99} ... \end{thebibliography} block (no external .bib file or
% bibtex pass needed -- \cite{} resolves against \bibitem entries in the same document via
% the normal two-pass pdflatex + .aux mechanism already used for every other cross-reference
% here). A lecture only lists the \bibitem entries it actually \cite{}s.
%
% To keep citation keys consistent as more lectures are added, reuse these canonical keys
% (and the same canonical \bibitem text) rather than inventing new ones for the same source:
%   db        Aho, Lam, Sethi, Ullman -- Compilers: Principles, Techniques, and Tools (2nd ed.)
%   et        Cooper & Torczon -- Engineering a Compiler (2nd ed.)
%   ap        Appel -- Modern Compiler Implementation in C/Java/ML
%   pp        Scott -- Programming Language Pragmatics
%   mu        Muchnick -- Advanced Compiler Design and Implementation
%   an        Parr -- The Definitive ANTLR 4 Reference
%   gr        Gries -- Compilers: A Survey
%   llvmlangref   LLVM Language Reference Manual
%   llvmka        "My First Language Frontend with LLVM" (Kaleidoscope Tutorial)
%   antlrsite     ANTLR4 project site
%   tssite        Tree-sitter site
%   mlir          Lattner et al., "MLIR: A Compiler Infrastructure for the End of Moore's Law"
%   mlirdocs      The MLIR Project, official documentation
% A "Citation ... multiply defined" warning in the combined full-course-notes.pdf build (when
% two lectures both cite, say, \bibitem[DB]{db}) is expected and harmless: each lecture's own
% thebibliography block is independent, the canonical text is identical everywhere it is
% reused, and the explicit [DB]-style label (not an auto-generated number) is what actually
% gets printed, so it renders identically regardless of which instance "wins" in the aux file.
